\section{Experimental Protocol}
\label{sec:setup}

\subsection{Designs, training arms, and sampling}

The benchmark covers streaming finite-impulse-response (FIR), reverse-FIR,
polynomial, recursive-filter, and median-filter accelerators under a uniform
clocked interface.  The families cover feed-forward arithmetic pipelines,
recurrences, and comparator networks.  Each specification and its executable
reference are generated programmatically, and several implementation styles
provide physically distinct but behaviorally equivalent training examples.
This construction lets us compare several accelerator structures within a
controlled domain, while keeping the scope distinct from arbitrary industrial RTL.

The sealed evaluation contains \studyclaim{StudyTwoDesignCount} designs,
balanced between \studyclaim{StudyTwoInterpDesignCount} interpolation and
\studyclaim{StudyTwoExtrapDesignCount} extrapolation tasks.  Family quotas are
unequal because the eligible median interpolation pool was empty; the median
family is represented only under extrapolation, so median-family and regime
effects cannot be separated.  The family breakdown in
Sec.~\ref{sec:results} reports the exact composition.  Interpolation changes parameters within
training ranges; extrapolation uses larger or newly parameterized instances
outside those ranges.  Every sealed design is excluded from the SFT corpus,
reward rows, policy prompts, and prior analysis.  The split remained unopened
until all declared checkpoints and pre-open validation artifacts were frozen.


The primary comparison starts from RTLCoder and its completion-masked SFT
adapter, and evaluates
two independently trained repaired-RF policy adapters. The MLP arm trains two
matched policies from the same SFT adapter using the original reward.  The
correctness-only arm uses one preregistered full-budget seed.  Arms use the same non-flat training-update target; flat groups provide no
policy-ranking information and do not advance the counter.  A crashed run is discarded and
restarted from the frozen SFT state; partial histories are never merged.
Equal update targets do not imply equal training compute: flat groups still
require generation and oracle execution, particularly for correctness-only
rewards.  The primary comparison matches deployment draws, not end-to-end
training time or energy.

At the sealed endpoint, SFT and every trained arm receive
\studyclaim{StudyTwoSamplesPerArmDesign} draws per design at temperature
\studyclaim{StudyTwoEvaluationTemperature}.  Two-seed arms contribute
\studyclaim{StudyTwoSamplesPerSeedDesign} draws per seed.  The midpoint
diagnostic uses \studyclaim{StudyTwoMidSamplesPerSeedDesign} draws per RF seed
and design.  Generation seeds are distinct from training seeds and from the
\studyclaim{StudyTwoOracleStreams} oracle streams, each of which contains
\studyclaim{StudyTwoOracleVectors} test vectors.  Accepted candidates are
deduplicated only for implementation; multiplicity is preserved for every
policy statistic.

\subsection{Post-route evaluation and statistical analysis}

We use the frozen physical flow \studyclaim{StudyTwoVivadoEnvironment}. Every
distinct correct candidate is synthesized, placed, and routed with the same
\claim{TargetPeriod}-ns requested clock.  Achievable frequency is derived from
the request and reported worst negative slack.  We therefore use
``post-route timing-derived \fmax'' rather than ``timing closure'': the flow
does not repeatedly rerun implementation at candidate-specific clock periods.
The historical script introduces the clock constraint after synthesis; this
flow is held fixed rather than retroactively corrected in the sealed result.

The primary estimand is Eq.~\eqref{eq:equal-sample}.  Correctness and
conditional frequency are diagnostics, not replacement endpoints.  The
primary RF-minus-SFT confidence plan is frozen and stratified by family and
regime, with percentile bounds.  RF-minus-MLP intervals reuse its exact
resample plans as a bounded post-primary comparison.  The family cells are
small and unequal, so the manuscript adds an unstratified paired-bootstrap and
leave-one-family-out sensitivity but does not replace the frozen primary plan.
No multiplicity adjustment is applied: the RF-versus-SFT contrast is the sole
prespecified primary inference, while control, family, PPA, and mechanism
analyses are secondary or descriptive.  Resource and power analyses condition
on designs with compiled correct support in SFT, MLP, and RF; assigning zero
area or power to a failed sample would incorrectly make failure look efficient.
Power is Vivado's vectorless estimate, not measured board power.


\subsection{Physical validation and additional seeds}

Three supporting studies probe physical sensitivity. A PYNQ-Z2 experiment
uses paired SFT/RF candidates selected by multiplicity and an outcome-independent
hash. All DUTs and an echo canary share one image and clock grid. Eligibility
and sweep coverage were amended before board measurement, so the subset is
descriptive. The Supplementary Material records the harness and amendments.
A separate pilot brackets setup-passing clock periods for
\revisionclaim{ClosureCandidateCount} selected candidates under Vivado
\revisionclaim{ClosureVersion}, with constraints applied before synthesis.
An exhaustive cross-version diagnostic retains the original RTL, device,
clock request, and constraint order, and checks setup-path coverage separately.

The additional-seed study evaluates RF seeds three and four and correctness-only
seeds two and three on all \revisionclaim{SeedDesignCount} frozen designs.
SFT is remeasured under the same bounded setup-period flow, retaining each
checkpoint's original draws and all physical failure zeros. The physical
packages were frozen before measurement; the statistical specification followed
inspection of point estimates. Its exploratory paired design-bootstrap
intervals use \revisionclaim{SeedBootstrapCount} shared family-by-regime
resamples. They condition on the sampled checkpoints and are reported separately
from the original timing-derived endpoint. None of these studies establishes
full hold and pulse-width timing sign-off.

\subsection{Transfer and supporting evidence}

To assess transfer beyond the main backbone, we evaluate the complete
SFT-to-policy recipe on
Qwen-Coder~\cite{hui2024qwen25coder} using the
earlier \claim{QwenDesignCount}-design, \claim{QwenFamilyCount}-family split.
This second-backbone experiment uses an earlier reward/training configuration,
so it tests transfer of the recipe rather than replication of the sealed RF
package. A matched VerilogEval extension tests the sealed SFT and RF checkpoints
on \revisionclaim{VerilogProblems} official problems using
\revisionclaim{VerilogSamples} draws per checkpoint and problem. The
Supplementary Material preserves the earlier five-family evaluations,
predictor diagnostics, API and HLS context, and detailed protocol history.
Those studies have distinct sampling or selection rules and are not pooled
with the sealed estimate. All reported numerical values are generated from
retained artifacts; complete denominators, failures, source hashes, and claim
provenance are checked during manuscript regeneration.
